\documentclass[11pt,a4paper]{article}

% ── Packages ─────────────────────────────────────────────────────────────────
\usepackage[utf8]{inputenc}
\usepackage[T1]{fontenc}
\usepackage[french]{babel}
\usepackage{geometry}
\geometry{top=2.5cm,bottom=2.5cm,left=2.5cm,right=2.5cm}

\usepackage{lmodern}
\usepackage{microtype}
\usepackage{graphicx}
\usepackage{booktabs}
\usepackage{tabularx}
\usepackage{array}
\usepackage{xcolor}
\usepackage{fancyhdr}
\usepackage{titlesec}
\usepackage{caption}
\usepackage{subcaption}
\usepackage{float}
\usepackage{amsmath}
\usepackage{amssymb}
\usepackage{hyperref}
\usepackage{enumitem}
\usepackage{tcolorbox}
\usepackage{mdframed}
\usepackage{parskip}

% ── Couleurs ──────────────────────────────────────────────────────────────────
\definecolor{meteoBlue}{HTML}{1D3557}
\definecolor{meteoRed}{HTML}{E63946}
\definecolor{meteoGray}{HTML}{F1FAEE}
\definecolor{meteoCyan}{HTML}{457B9D}

% ── Mise en page ─────────────────────────────────────────────────────────────
\pagestyle{fancy}
\fancyhf{}
\fancyhead[L]{\small\color{meteoBlue}\textbf{Météorage} — Prédiction de fin d'orage}
\fancyhead[R]{\small\color{meteoBlue}Analyse exploratoire des données}
\fancyfoot[C]{\small\color{gray}\thepage}
\renewcommand{\headrulewidth}{0.4pt}
\renewcommand{\footrulewidth}{0pt}

% ── Titres ────────────────────────────────────────────────────────────────────
\titleformat{\section}[block]
  {\Large\bfseries\color{meteoBlue}}
  {\thesection}{1em}{}[\vspace{-0.3em}\rule{\linewidth}{0.4pt}]

\titleformat{\subsection}[block]
  {\large\bfseries\color{meteoCyan}}
  {\thesubsection}{1em}{}

\titleformat{\subsubsection}[block]
  {\normalsize\bfseries\color{meteoBlue}}
  {\thesubsubsection}{1em}{}

% ── Légendes ──────────────────────────────────────────────────────────────────
\captionsetup{font=small, labelfont=bf, margin=10pt}

% ── Boîte de synthèse ─────────────────────────────────────────────────────────
\newenvironment{keybox}[1]
  {\begin{tcolorbox}[colback=meteoGray,colframe=meteoBlue,title={\textbf{#1}},
    fonttitle=\small, boxrule=0.6pt, arc=3pt]}
  {\end{tcolorbox}}

% ─────────────────────────────────────────────────────────────────────────────
\begin{document}

% ── Page de titre ─────────────────────────────────────────────────────────────
\begin{titlepage}
  \pagecolor{meteoBlue}
  \color{white}
  \vspace*{2cm}
  \begin{center}
    {\Huge\bfseries Prédiction de fin d'orage\\[0.4em]}
    {\large\bfseries Rapport d'Analyse Exploratoire des Données}\\[2em]
    \rule{0.5\linewidth}{0.8pt}\\[1.5em]
    {\large Réponse à l'appel d'offre \textbf{Météorage}}\\[0.5em]
    {\normalsize Analyse de la dynamique spatio-temporelle des éclairs\\
    autour de 5 aéroports européens}\\[3em]
    {\normalsize Mars 2026}
  \end{center}
  \vfill
  \begin{center}
    \small\textit{Données : segment\_alerts\_all\_airports\_train.csv}\\
    \small\textit{507\,071 éclairs $\cdot$ 5 aéroports $\cdot$ 769 sessions d'alerte}
  \end{center}
\end{titlepage}
\nopagecolor

% ── Table des matières ────────────────────────────────────────────────────────
\tableofcontents
\newpage

% =============================================================================
\section{Introduction et contexte}
% =============================================================================

Le problème posé par Météorage est celui de l'\textbf{estimation probabiliste de la fin d'un orage}. 
Aujourd'hui, les alertes restent actives une durée fixe (30 à 60 minutes) après le dernier éclair 
cloud-to-ground (\textit{CG}) dans la zone de surveillance. Pour des secteurs critiques comme les aéroports, 
cette approche conservatrice engendre des pertes d'exploitation significatives.

L'objectif est de développer un modèle capable d'exploiter la \textbf{dynamique spatio-temporelle} 
des éclairs pour estimer, en temps réel, la probabilité que l'orage soit terminé.

\subsection{Description du jeu de données}

\begin{table}[H]
\centering
\renewcommand{\arraystretch}{1.3}
\begin{tabularx}{\linewidth}{>{\bfseries}lXl}
\toprule
Variable & Description & Type \\
\midrule
\texttt{lightning\_id} & Identifiant unique global de l'éclair & int \\
\texttt{lightning\_airport\_id} & Identifiant éclair dans le contexte aéroport & int \\
\texttt{date} & Timestamp UTC de l'éclair & datetime \\
\texttt{lon}, \texttt{lat} & Coordonnées géographiques de l'éclair & float \\
\texttt{amplitude} & Courant de foudre en kA (signe = polarité) & float \\
\texttt{maxis} & Amplitude normalisée / indicateur d'intensité & float \\
\texttt{icloud} & \texttt{True} = intra-cloud, \texttt{False} = cloud-to-ground & bool \\
\texttt{dist} & Distance à l'aéroport (km, rayon max 50 km) & float \\
\texttt{azimuth} & Angle par rapport au Nord (\degre) & float \\
\texttt{airport} & Nom de l'aéroport (5 valeurs) & str \\
\texttt{airport\_alert\_id} & Identifiant de la session d'alerte (NaN si hors alerte) & float \\
\texttt{is\_last\_lightning\_cloud\_ground} & \textbf{Cible} : dernier éclair CG de la session ? & bool/NaN \\
\bottomrule
\end{tabularx}
\caption{Description des variables du jeu de données}
\end{table}

\begin{keybox}{Statistiques globales}
\begin{itemize}[noitemsep]
  \item \textbf{507\,071} éclairs au total
  \item \textbf{5 aéroports} : Ajaccio, Bastia, Biarritz, Nantes, Pise
  \item \textbf{769 sessions d'alerte} distinctes (\texttt{airport\_alert\_id} non nul)
  \item \textbf{56\,599 éclairs en alerte} ($\approx 11\%$ du total)
  \item La variable cible est disponible uniquement pour les éclairs CG en session d'alerte
  \item Période couverte : 2016 (au moins) avec pics estivaux marqués
\end{itemize}
\end{keybox}

% =============================================================================
\section{Répartition par aéroport}
% =============================================================================

\begin{figure}[H]
  \centering
  \includegraphics[width=\linewidth]{figures/fig1_airport_counts.pdf}
  \caption{Distribution des éclairs entre les 5 aéroports. Pise concentre la majorité de l'activité 
  orageuse.}
\end{figure}

Les 5 aéroports présentent des niveaux d'activité orageuse très différents. Pise (Toscane, Italie) 
domine avec plus de 30\% des éclairs totaux, en raison de sa position géographique favorable au 
développement convectif (convergence mer Ligure / relief apennin). Bastia et Biarritz suivent avec 
des niveaux d'activité similaires, puis Ajaccio, et enfin Nantes, qui connaît l'activité la plus faible.
On remarque que Bastia et Ajaccio, les deux aéroports corses, ont des niveaux d'activité orageuse 
différentes (presque 2 fois plus d'éclairs à Bastia qu'à Ajaccio), bien qu'on verra plus tard que le
profil temporel est très similaire. 

% =============================================================================
\section{Analyse temporelle}
% =============================================================================

\subsection{Distribution mensuelle et horaire}

\begin{figure}[H]
  \centering
  \includegraphics[width=\linewidth]{figures/fig2_temporal.pdf}
  \caption{Distribution temporelle des éclairs : saisonnalité marquée avec un pic en été 
  (juin-septembre) et un cycle diurne prononcé avec maximum en fin d'après-midi.}
\end{figure}

Deux patterns temporels fondamentaux se dégagent :

\begin{itemize}
  \item \textbf{Saisonnalité estivale} : la convection thermique atteint son maximum en été, 
  avec un pic en juillet-août. La période novembre-mars est quasi inerte.
  \item \textbf{Cycle diurne convectif} : le maximum d'éclairs se situe dans l'après-midi, 
  entre 12h et 14h30 UTC, correspondant au pic de la couche limite convective après son réchauffement.
\end{itemize}

\subsection{Heatmap mois $\times$ heure}

\begin{figure}[H]
  \centering
  \includegraphics[width=\linewidth]{figures/fig3_heatmap_month_hour.pdf}
  \caption{Densité combinée mois * heure. Le quadrant "été / après-midi" concentre l'essentiel 
  de l'activité. Les orages nocturnes (00h-06h UTC) sont plus fréquents en juillet-août, 
  associés aux systèmes convectifs méso-échelle.}
\end{figure}

\subsection{Activité par aéroport et par mois}

\begin{figure}[H]
  \centering
  \includegraphics[width=\linewidth]{figures/fig4_heatmap_airport_month.pdf}
  \caption{Répartition mensuelle de l'activité orageuse par aéroport. On observe que Biarritz 
  présente un pic précoce et très étalé s'amorçant dès mai-juin jusqu'en août. Nantes a également 
  une activité concentrée sur le début d'été (mai--juillet). Pise connaît un pic tardif en août--septembre, 
  tandis qu'Ajaccio et Bastia, avec des profils quasi-identiques, voient l'essentiel de leur 
  activité orageuse repoussé de août à octobre, période de fort flux de chaleur de la mer Méditerranée.}
\end{figure}

% =============================================================================
\section{Caractéristiques physiques des éclairs}
% =============================================================================

\begin{figure}[H]
  \centering
  \includegraphics[width=\linewidth]{figures/fig5_physics.pdf}
  \caption{Distributions des variables physiques. \textbf{Amplitude} : distribution bimodale 
  asymétrique, la majorité des éclairs CG sont négatifs. \textbf{Maxis} : variable interne concentrée entre 0 et 2 avec une queue lourde. 
  \textbf{Distance} : distribution croissante linéairement — ce phénomène est un artefact géométrique attendu car 
  la surface de l'anneau de distance $r$ est proportionnelle à $2\pi r$. Pour une densité spatiale d'éclairs parfaitement 
  homogène, le nombre d'éclairs augmente linéairement avec la distance. 
  \textbf{Ratio IC} : varie de 66\% (Ajaccio) à plus de 80\% (Biarritz).}
\end{figure}

\subsection{Définition de Maxis et Ratio IC}

\textbf{Maxis} est un indicateur de gradient d'intensité couramment utilisé par Météorage. Il s'agit 
d'une valeur normalisée (souvent corrélée à l'amplitude maximale ou à la raideur du front de courant 
de l'éclair) qui aide à distinguer différents types ou puissances de décharges.

Le rapport éclairs intra-cloud (\textbf{IC}) sur le total d'éclairs est un indicateur du stade de développement 
d'un orage. Un ratio élevé en début de session indique un orage en phase de maturation ; 
un ratio croissant en fin de session peut signaler la dissipation.

\subsection{Rose des azimuts}

\begin{figure}[H]
  \centering
  \includegraphics[width=\linewidth]{figures/fig6_azimuth_rose.pdf}
  \caption{Distribution angulaire (azimut) des éclairs. Nantes affiche un fort flux en provenance du sud-ouest. 
  Biarritz se distingue par un flux d'est dominant (descente des Pyrénées). Pour la Corse, Bastia 
  présente des apports sud-ouest et sud-est, tandis qu'Ajaccio est majoritairement à l'est. Pise 
  connaît une répartition plus symétrique avec deux axes dominants est/ouest.}
\end{figure}

% =============================================================================
\section{Analyse géographique}
% =============================================================================

\begin{figure}[H]
  \centering
  \includegraphics[width=\linewidth]{figures/fig7_geo_scatter.pdf}
  \caption{Densité spatiale des éclairs dans un rayon d'environ 35 km autour de chaque aéroport. 
  L'étoile rouge indique la position de l'aéroport. On note des clusters intenses liés au relief (Corse, Apennins).}
\end{figure}

Les cartes révèlent des structures spatiales propres à chaque site : 
Ajaccio et Bastia montrent des éclairs concentrés sur les reliefs intérieurs corses ; 
Pise présente une dispersion vers les Apennins à l'est ; Biarritz et Nantes ont 
des distributions plus homogènes, typiques des environnements de plaine atlantique.

% =============================================================================
\newpage
\section{Analyse des sessions d'alerte}
% =============================================================================

\subsection{Durée et intensité des sessions}

\begin{figure}[H]
  \centering
  \includegraphics[width=\linewidth]{figures/fig8_sessions.pdf}
  \caption{Caractéristiques des 769 sessions d'alerte. Le nombre d'éclairs par session suit 
  une forte loi asymétrique, avec de très nombreuses sessions mineures, mais aussi des sessions majeures 
  atteignant un millier d'éclairs.}
\end{figure}

\begin{figure}[H]
  \centering
  \includegraphics[width=\linewidth]{figures/fig9_sessions_by_airport.pdf}
  \caption{Comparaison des sessions par aéroport. Pise présente les sessions les plus 
  longues et les plus intenses (médiane élevée en nombre d'éclairs). Biarritz et Nantes 
  ont des sessions plus courtes mais potentiellement plus imprévisibles.}
\end{figure}

\subsection{Spécificités par aéroport}

\begin{keybox}{Profils orageaux par aéroport}
\begin{itemize}[noitemsep]
  \item \textbf{Pise} : L'aéroport le plus actif en volume absolu, avec un pic de fin d'été.
  \item \textbf{Ajaccio \& Bastia} : Saisonnalité quasi-identique (août à octobre).
  \item \textbf{Nantes} : Orages plus rares mais concentrés en début d'été (mai-juillet), marqués 
  par un flux SW puissant.
  \item \textbf{Biarritz} : Fort ratio de décharges IC et vents orageux d'est (effet pyrénéen 
  et transition océan/continent).
\end{itemize}
\end{keybox}

% =============================================================================
\section{Analyse de la variable cible}
% =============================================================================

\begin{figure}[H]
  \centering
  \includegraphics[width=\linewidth]{figures/fig10_target_analysis.pdf}
  \caption{\textbf{Gauche} : la cible est fortement déséquilibrée — le dernier éclair CG 
  ne représente qu'une fraction des éclairs en session (par définition, un par session). 
  \textbf{Centre} : distribution du temps restant jusqu'à la fin de session — 
  la règle fixe d'attente de 30 minutes après un éclair ne capture que 25.0\% du temps restant total de la tempête ; une prédiction dynamique 
  peut réduire significativement ce décalage. 
  \textbf{Droite} : intervalles inter-éclairs CG, avec une queue lourde significative.}
\end{figure}

\subsection{Formulation du problème de modélisation}

Deux formulations complémentaires sont envisagées pour la Phase 2 :

\begin{enumerate}
  \item \textbf{Classification binaire} : pour chaque éclair CG en session, 
  prédire $P(\text{dernier éclair CG}) \in [0,1]$. 
  Le modèle produit une probabilité de fin d'orage en temps réel.
  
  \item \textbf{Régression de survie} : modéliser le temps restant $T$ jusqu'au 
  dernier éclair CG comme une variable de survie censurée, avec :
  \[
    P(T > t \mid \mathcal{F}_t) = S(t \mid \mathcal{F}_t)
  \]
  où $\mathcal{F}_t$ est l'historique des éclairs jusqu'au temps $t$.
\end{enumerate}

\subsection{Exemple de session d'alerte}

\begin{figure}[H]
  \centering
  \includegraphics[width=\linewidth]{figures/fig11_session_example.pdf}
  \caption{Exemple d'une session d'alerte complète. Le panneau supérieur montre 
  la distance à l'aéroport des éclairs IC (violet) et CG (rouge) au cours du temps. 
  Le panneau inférieur montre l'évolution de l'amplitude des éclairs CG. 
  La ligne pointillée noire indique le moment du dernier éclair CG — 
  c'est ce que le modèle doit apprendre à anticiper.}
\end{figure}

% =============================================================================
\section{Corrélations et structure des données}
% =============================================================================

\begin{figure}[H]
  \centering
  \includegraphics[width=0.6\linewidth]{figures/fig12_correlation.pdf}
  \caption{Matrice de corrélation de Pearson des variables physiques. Contrairement à l'attente initiale, 
  l'amplitude et \texttt{maxis} présentent une corrélation quasi-inexistante ($0.08$). Cela 
  signifie qu'elles mesurent des caractéristiques très distinctes de la décharge (puissance vs profil). 
  La distance et l'azimuth sont également quasi-indépendants des variables physiques.}
\end{figure}

% =============================================================================
\section{Conclusions et orientations pour la modélisation}
% =============================================================================

\begin{keybox}{Points clés de l'analyse}
\begin{itemize}[noitemsep]
  \item \textbf{Forte saisonnalité} : les modèles devront être entraînés/évalués 
  en respectant la structure temporelle (pas de leakage temporel dans la validation).
  \item \textbf{Hétérogénéité inter-aéroports} : un encodage de l'aéroport ou des 
  modèles spécifiques par site seront nécessaires.
  \item \textbf{Déséquilibre de classes} : la cible est déséquilibrée par construction 
  (1 dernier éclair par session). Des techniques de pondération ou d'échantillonnage 
  seront requises.
  \item \textbf{Structure séquentielle} : l'ordre temporel des éclairs dans une session 
  est l'information la plus riche. Les modèles devront intégrer cette dimension 
  (features de fenêtre glissante, processus ponctuels, RNN/Transformer...).
  \item \textbf{Signal de dissipation} : le ratio IC croissant en fin de session 
  et l'augmentation des intervalles inter-éclairs CG sont des signaux prometteurs.
\end{itemize}
\end{keybox}

\subsection{Features candidates pour la modélisation}

\begin{table}[H]
\centering
\renewcommand{\arraystretch}{1.3}
\begin{tabularx}{\linewidth}{>{\ttfamily\small}lX}
\toprule
\textbf{Feature} & \textbf{Description} \\
\midrule
time\_since\_alert\_start & Temps écoulé depuis le début de la session (min) \\
time\_since\_last\_cg & Intervalle depuis le dernier éclair CG (min) \\
n\_cg\_last\_10min & Nombre d'éclairs CG dans les 10 dernières minutes \\
n\_ic\_last\_10min & Nombre d'éclairs IC dans les 10 dernières minutes \\
icloud\_ratio\_recent & Ratio IC/(IC+CG) sur fenêtre glissante \\
mean\_dist\_recent & Distance moyenne des éclairs récents \\
dist\_centroid\_drift & Déplacement du centroïde des éclairs récents \\
amplitude\_trend & Tendance (pente) de l'amplitude sur les derniers CG \\
maxis\_trend & Tendance de maxis \\
airport & Encodage de l'aéroport \\
hour\_utc & Heure UTC (signal diurne) \\
month & Mois (signal saisonnier) \\
\bottomrule
\end{tabularx}
\caption{Features candidates pour les modèles de Phase 2}
\end{table}

% =============================================================================
\appendix
\section{Environnement technique}
% =============================================================================

\begin{itemize}
  \item Langage : Python 3.14, environnement \texttt{pew orage}
  \item Librairies : \texttt{pandas 3.0}, \texttt{numpy}, \texttt{matplotlib}, 
    \texttt{seaborn}, \texttt{scikit-learn}, \texttt{scipy}
  \item Compilation : MiKTeX 26.2 (pdfLaTeX)
\end{itemize}

\end{document}
